\chapter{氛围动力学：社会信号、环境线索与适应轨迹}
\label{chp:atmosphere_geometric_dynamics}
上一章建立了低活动、高增益与现实校准的条件，本章把视线从单个系统移向人与人、系统与空间的连续耦合。日常所谓“气场”“氛围”“压迫感”并非无通道传播的物理场，也不能被轻率贬为纯粹错觉；它们可以来自姿态、语音、距离、建筑、制度符号、交互历史和接收者预期共同塑造的条件化体验。本章以计算动力学重建四类问题：社会信号如何跨接口传递，环境线索如何改变行动分布，人际协调为何产生不同负担，以及长期暴露如何与可塑性、发展阶段和保护因素共同形成适应轨迹，并给出可反驳的边界。

\section{气场的可观测发生学：社会信号、通道与接收者更新}
\label{sec:physics_of_aura}
我们首先取消“规范场非定域辐射”的字面主张。发送者内部目标、情绪和信念属于其自身状态，接收者只有通过可用通道取得观测，才可能更新对发送者的估计。通道包括语言内容、声学韵律、注视、姿态、动作时序、空间距离、服饰与身份标记，也包括第三方记录和共同情境；它们都受传播速度、采样精度、遮挡、文化规则和主动伪装限制。因此，气场不是一种脱离载体的能量，而是接收者对一束多模态社会信号及其不确定性的整体读出。

设发送者A在时刻 $k$ 的内部状态为 $x_k^A$，可执行动作为 $a_k^A$，物理与数字通道为 $C_k$，接收者B取得的观测为
\[
y_k^B=C_k\,g_A(x_k^A,a_k^A)+n_k,
\]
其中 $g_A$ 是内部状态到可观察行为的生成映射，$n_k$ 汇集环境噪声、设备误差与遮挡。$x_k^A$ 不可被B直接访问，且同一行为可能由不同意图产生；微笑可能表示友好、礼貌、紧张或策略性掩饰。接收者只能维护后验 $p_B(x_k^A\mid y_{0:k}^B,c_k)$，其中 $c_k$ 是场景、关系史和文化语境。气场体验对应这个后验对B的注意、价值评估和行动门控产生的联合影响，而不是A把自己的内部几何复制到B。

原章把坚定信念、低内耗和精确行为称为低熵应力张量。这里要分开三个量：行为序列的统计可预测性，表达与目标的一致性，以及执行所消耗的实际资源。一个人语速稳定、动作节律清楚、言行一致，可能使观察者更容易预测其下一步，从而形成“笃定”的印象；但低统计熵不自动意味着正确、善意或低能耗，训练有素的欺骗者也能给出高度一致的信号。可靠判断必须比较多时段、多情境和受约束行为，检查表达在反例、利益冲突和承诺成本下是否保持。

接收端更新可写为
\[
z_{k+1}^B=F_v^B(z_k^B,G_k^B,y_k^B,c_k)+\xi_k^B,
\qquad r_k^B=R^B(z_k^B,G_k^B),
\]
其中 $z_k^B$ 是全局分布式激活，$G_k^B$ 是保存人物身份、关系史和承诺的局部结构网络，$r_k^B$ 是敬畏、亲近、警惕或回避等任务相关读出。两类表示通过人物身份和时间戳绑定：分布表示提供快速相似判断，结构网络防止把不同人的表情、旧事件和当前语句混为一谈。这个更新也解释了为何同一发送者会让不同接收者产生不同体验；B的先验、创伤史、权力关系、目标和当下资源都会改变 $F_v^B$，因此不存在脱离接收者的绝对气场强度。

原章区分同相共振、拓扑压制与非对易冲突，我们把它们改写为三类可测响应。第一类是协调：信号与B当前目标、规范和预测兼容，理解时延和冲突修复成本下降，B可能感到被鼓舞或轻松。第二类是权力约束：A掌握奖惩、信息或行动入口，B即使不认同也会收缩可选动作，体验为敬畏或压迫。第三类是规范冲突：双方目标、承诺或角色预期不兼容，B必须投入监测、反驳和边界维护，主观负担上升。三类可以同时存在，例如一位导师既提供高质量协调，也拥有评价权；不能仅凭“舒服”判断关系有益，也不能把紧张都归因于价值正交。

为检验气场效应，我们至少需要发送者、通道和接收者三侧的测量。发送者侧记录可观察行为与情境，不把自述当成唯一真值；通道侧测量信号强度、延迟、缺失和压缩；接收者侧记录判断、置信度、生理或行为变化以及实际决策。实验可用遮蔽面孔、去除语义、改变身份提示、随机播放顺序或交换房间来做消融。如果只保留声音后效应消失，说明视觉或身份线索重要；若告知不同头衔就改变评价，权威先验可能强于行为本身。跨文化重复还能检查所谓普遍信号是否只是本地规范。

在AgentOS实例中，SPC维护发送者身份和会话关系，TCE把当前目标、风险与权力约束纳入评估，CCM控制多模态候选和反解释的并发，Boundary限制在低置信社会判断下可执行的动作，Offloading把原始观测、推断和后续回执分层写入外部记忆 $\mathrm{Me}$。$h(t)$ 承载即时交互态，$\Psi$ 保存可撤销候选，$\Theta$ 只在证据跨时段一致时提交关系结构编辑。这是一种工程实现，不是所有社会智能的唯一机制。

因此，气场的“相对性”不是场论方程的证明，而是一项识别原则：观察到的影响由发送行为、传输通道、接收模型和情境共同决定。反例包括无人在场时仍由预先信息产生的紧张、同一录音配不同身份标签引发的不同判断，以及强烈但错误的第一印象。把这些反例纳入模型后，我们既不需要诉诸超距意念，也不必否认整体体验的真实性；需要拒绝的是从体验强度直接跳到发送者本质、道德品质或不可见物理场的推断。
